\documentclass[11pt]{article}
\usepackage[margin=1in]{geometry}
\usepackage{amsmath,amssymb,amsthm}
\usepackage{booktabs}
\usepackage{graphicx}
\usepackage[numbers,sort&compress]{natbib}
\usepackage{hyperref}
\usepackage{xcolor}
\usepackage{algorithm}
\usepackage{algpseudocode}
\usepackage{adjustbox}

\newtheorem{theorem}{Theorem}
\newtheorem{proposition}[theorem]{Proposition}
\newtheorem{corollary}[theorem]{Corollary}
\newtheorem{assumption}{Assumption}
\theoremstyle{definition}
\newtheorem{definition}{Definition}
\newtheorem{remark}{Remark}

\newcommand{\E}{\mathbb{E}}
\newcommand{\Var}{\operatorname{Var}}
\newcommand{\diag}{\operatorname{diag}}
\newcommand{\tW}{\tilde W}
\newcommand{\tY}{\tilde Y}

\title{Walk-Unbiased Doubly Robust Graph Propagation:\\
Provably Fair, Transparent, Accountable and Private Recommendation from Exposure-Biased Logs}
\author{Quang-Vinh Dang\\
British University Vietnam, Hung Yen, Vietnam\\
\texttt{vinh.dq4@buv.edu.vn} \quad ORCID: 0000-0002-3877-8024}
\date{}

\begin{document}
\maketitle

\begin{abstract}
Graph-based recommenders propagate preferences along multi-hop walks of the
logged user--item graph. This graph is filtered by the platform's own exposure
policy, so propagation amplifies exposure bias. Recent debiasing methods
reweight the edges of this graph by inverse propensities or impute missing
edges with doubly robust (DR) estimates. We show that these corrections are
still biased once propagation spans more than one hop. A walk that revisits an
edge squares that edge's estimate, and the squared inverse propensity inflates
exactly the rarely exposed edges that debiasing is meant to recover. We derive
the exact bias and a lower bound that grows as the inverse of the propensity
clip. Because outcomes are binary, the fix is simple: every repeated edge is
replaced by its unbiased idempotent counterpart. The resulting operator,
\emph{DRUP} (Doubly Robust, walk-Unbiased Propagation), is exactly unbiased and
edge-wise doubly robust for propagation over the counterfactual full-exposure
graph, and it costs no more than the uncorrected operator. The same structure
gives DRUP provable properties beyond accuracy:
(i) variance, clipping-bias and concentration bounds;
(ii) exposure invariance, a causal item-fairness guarantee that the expected
score does not depend on interventions on an item's exposure, whereas logged-graph
propagation amplifies exposure and uncorrected DR propagation over-corrects it;
(iii) scores that are affine in each user's own interactions, which yields exact
Shapley attributions and provably minimal counterfactual explanations;
(iv) closed-form certificates against injected fake profiles; and
(v) a single public item operator whose Gaussian release makes the whole
recommender jointly differentially private.
On Coat and KuaiRec, whose test sets are random or fully observed, DRUP needs no
training. It is accuracy-competitive with trained MF, LightGCN, NAVIP and DR
baselines (best on KuaiRec, tied with the best trained model on Coat), and
every guarantee is verified empirically.
\end{abstract}

\noindent\textbf{Keywords:} causal recommendation; graph neural networks; doubly robust estimation;
exposure bias; fairness; explainability; certified robustness; differential privacy.

%======================================================================
\section{Introduction}
%======================================================================

Graph-based recommender systems (GBRs) represent users and items as nodes and
logged interactions as edges. They score a user--item pair by aggregating
signals along multi-hop walks~\cite{he2020lightgcn,shen2021gfcf}. The logged
graph, however, is not the preference graph. An edge exists only if the platform
\emph{exposed} the item and the user then responded. A recent survey of causal
learning in GBRs~\cite{yu2026survey} traces four open challenges to this gap:
out-of-distribution generalisation, dynamic preferences, fairness and
interpretability. All four arise because message passing propagates
correlations of the observational graph rather than causal preference
structure. The same survey lists the lack of theoretical guarantees,
scalability and privacy as open directions.

Two families of fixes exist. \emph{Loss-level} debiasing (IPS, DR) reweights
the training objective~\cite{schnabel2016,wang2019drjl,saito2020unbiased} but
leaves the propagation operator biased. \emph{Propagation-level} debiasing
reweights edges by inverse propensities during neighbour
aggregation~\cite{kim2022navip,dpaa2026}. We show that the second family is
itself biased for multi-hop propagation. The root cause is elementary. The
edge estimate $W_e=O_eY_e/p_e$ is unbiased for the potential outcome $Y_e$, but
a walk that traverses $e$ twice contributes $W_e^2$, and
$\E[W_e^2]=Y_e/p_e\neq Y_e$. Walks with repeated edges are not rare corner
cases. They include every back-and-forth step $u\to j\to u$ that makes LightGCN
self-reinforcing. Their bias is largest precisely where the propensity is
smallest.

\paragraph{Contributions.}
\begin{enumerate}
\item We identify the \emph{repeated-walk bias} of inverse-propensity and DR
graph propagation. We give its exact form and a lower bound that is unbounded as
propensities vanish (Theorem~\ref{thm:lower}).
\item We propose DRUP. Using $Y_e^2=Y_e$, DRUP replaces each repeated edge by
its idempotent unbiased counterpart. The resulting 3-hop operator is exactly
unbiased and edge-wise doubly robust for propagation over the counterfactual
full-exposure graph (Theorem~\ref{thm:unbiased}). It has the cost of one item
Gram matrix, needs no training, and factorises into a public item operator plus
per-user local corrections.
\item We prove variance, clipping-bias and concentration bounds
(Theorem~\ref{thm:var}).
\item We show that the same structure yields provable \emph{FAT+P}
properties:
causal item fairness via exposure invariance (Theorem~\ref{thm:fair});
exact attributions and optimal counterfactual explanations
(Proposition~\ref{prop:affine}, Theorem~\ref{thm:cf});
certified robustness to injected profiles (Theorem~\ref{thm:robust});
and joint differential privacy (Theorem~\ref{thm:dp}).
\item We evaluate DRUP on Coat~\cite{schnabel2016}, whose test set is a random
sample, and on KuaiRec~\cite{gao2022kuairec}, whose test matrix is fully
observed. We compare it with trained and training-free baselines and verify each
guarantee: by Monte Carlo, by a real-data intervention on exposure, by brute
force, and against realised attacks.
\end{enumerate}

%======================================================================
\section{Related work}
%======================================================================

\paragraph{Causal and debiased recommendation.}
Inverse propensity scoring~\cite{schnabel2016} and doubly robust
learning~\cite{wang2019drjl} are the standard tools for MNAR feedback. Causal
embeddings disentangle interest from conformity~\cite{zheng2021dice}. The survey
of~\cite{yu2026survey} organises causal GBRs by challenge. None of these
methods analyses the bias of the \emph{propagation operator} beyond one hop.

\paragraph{Debiased message passing.}
NAVIP~\cite{kim2022navip} applies inverse-propensity weights with Laplacian
normalisation inside LightGCN aggregation. DPAA~\cite{dpaa2026} learns
popularity-aware aggregation weights. Both are unbiased for a single
aggregation step at most. Theorem~\ref{thm:lower} shows that stacking such
steps reintroduces a bias of order $1/p$.

\paragraph{Linear graph filters.}
LightGCN is a polynomial filter of the normalised adjacency, and closed-form
filters are competitive with it~\cite{shen2021gfcf}. DRUP is a closed-form
filter on a debiased operator. Its linearity is what makes exact explanations,
certificates and DP releases possible.

\paragraph{Explanations, robustness and privacy.}
GNNExplainer~\cite{ying2019gnnexplainer} and CF-GNNExplainer~\cite{lucic2022cfgnn}
approximate explanations by optimising on a surrogate. PORE~\cite{jia2023pore}
certifies recommenders against poisoning by ensembling. Differentially private
recommenders release noisy item covariances~\cite{mcsherry2009netflix}.
DRUP gives exact explanations, closed-form certificates and a joint-DP
guarantee~\cite{hsu2014matchings} for one and the same scoring rule.

%======================================================================
\section{Problem setting}
%======================================================================

Users $u\in[m]$, items $i\in[n]$ and pairs $e=(u,i)$ carry a binary potential
outcome $Y_e$ and an exposure $O_e\sim\mathrm{Bernoulli}(p_e)$. The log reveals
$O$ and $O\odot Y$.

\begin{assumption}[Unconfoundedness]\label{as:unconf}
Given the propensities, the exposures $O_e$ are mutually independent and
independent of $Y$.
\end{assumption}

\begin{assumption}[Fixed nuisances]\label{as:fixed}
The propensity estimate $\hat p$, the imputation $\hat Y\in[0,1]^{m\times n}$
and the edge weights $C$ do not depend on $O$. This holds with cross-fitting or
with nuisances estimated on a vetted historical log. We use
$C_{ui}=d_u^{-\alpha}d_i^{-(1-\alpha)}$ with degrees taken from $\hat Y$.
\end{assumption}

With $\tY=C\odot Y$, the target is propagation over the full-exposure graph:
\begin{equation}
F^*=a\,\tY+b\,T^*,\qquad T^*=\tY\tY^\top\tY .
\end{equation}
This is the linear LightGCN / graph-filter score one would compute if every
item had been shown to every user.

%======================================================================
\section{DRUP}
%======================================================================

\paragraph{Edge estimates.}
With clipped propensities $\bar p=\max(\hat p,\tau)$,
\begin{equation}
W^{\mathrm{IPS}}_e=\frac{O_eY_e}{\bar p_e},\qquad
W^{\mathrm{DR}}_e=\hat Y_e+\frac{O_e(Y_e-\hat Y_e)}{\bar p_e}.
\end{equation}
Both have mean $Y_e$ when $\bar p=p$.

\paragraph{Idempotent walk correction.}
Let $\tW=C\odot W$ and $\Delta=\tW\odot\tW-C\odot C\odot W$, with row sums $r$
and column sums $\kappa$. The corrected 3-hop operator is
\begin{equation}\label{eq:drup}
\hat T=\tW\tW^\top\tW-\tW\odot(r\mathbf 1^\top-\Delta)-\tW\odot(\mathbf 1\kappa^\top-\Delta)
-(\tW^{\odot3}-C^{\odot3}\odot W).
\end{equation}
The three corrections handle walks $u\!-\!j\!-\!v\!-\!i$ with $v=u$, with $j=i$,
and with both. In each case every repeated factor $\tW_e^k$ is replaced by
$C_e^kW_e$. The DRUP score is $s=a\tW+b\hat T$ with global constants $a,b$.

\paragraph{Public-operator form.}
Define the diagonal-corrected item Gram matrix
$G=\tW^\top\tW-\diag(\tW^\top\tW)+\diag(\mathbf 1^\top(C\odot C\odot W))$.
Then
\begin{equation}\label{eq:local}
\hat T_{u\cdot}=\tilde w_uG-\tilde w_u\odot(r_u\mathbf 1^\top-2\delta_u)-(\tilde w_u^{\odot3}-c_u^{\odot3}\odot w_u),
\end{equation}
where $\delta_u=\tilde w_u^{\odot2}-c_u^{\odot2}\odot w_u$. Equation~\eqref{eq:local}
matches~\eqref{eq:drup} to machine precision. $G$ is the only object shared
across users (Algorithm~\ref{alg:drup}).

\begin{algorithm}[t]
\caption{DRUP (training-free)}\label{alg:drup}
\begin{algorithmic}[1]
\Require log $(O,O\odot Y)$; hyper-parameters $\tau,\alpha,\lambda,b/a$
\State fit propensities $\hat p$ (logistic exposure model) and imputation $\hat Y$ (IPS-weighted additive ridge)
\State $\bar p\gets\max(\hat p,\tau)$; \ $W\gets\hat Y+O\odot(Y-\hat Y)/\bar p$
\State $C\gets d_u^{-\alpha}d_i^{-(1-\alpha)}$ with $d$ from $\hat Y$
\State $G\gets\tW^\top\tW-\diag(\tW^\top\tW)+\diag(\mathbf 1^\top(C\odot C\odot W))$ \Comment{public operator}
\For{each user $u$ (locally)}
\State $\hat T_{u\cdot}\gets$ Eq.~\eqref{eq:local}; \ $s_u\gets a\,\tilde w_u+b\,\hat T_{u\cdot}$; rank unexposed items by $s_u$
\EndFor
\end{algorithmic}
\end{algorithm}

\paragraph{Complexity.}
One $n\times n$ Gram product and $O(mn)$ reductions, i.e., the cost of the
uncorrected operator. On KuaiRec (7{,}176 users, 10{,}728 items, 10.3M exposures)
one configuration takes under a minute on four CPU cores.

%======================================================================
\section{Theory}
%======================================================================

Full proofs are in Appendix~\ref{app:proofs}. Let $\varepsilon_e=|Y_e-\hat Y_e|$.

\begin{theorem}[Edge-wise double robustness]\label{thm:unbiased}
Under Assumptions~\ref{as:unconf}--\ref{as:fixed}, let
$\delta_e=(p_e/\bar p_e-1)(Y_e-\hat Y_e)$. Then $\E[\hat T_{ui}]$ equals $T^*_{ui}$
with every $Y_e$ replaced by $Y_e+\delta_e$. In particular, if every edge has a
correct propensity or a correct imputation, then $\E[s]=F^*$.
\end{theorem}

\begin{theorem}[Repeated-walk bias of uncorrected propagation]\label{thm:lower}
With correct propensities and a candidate $(u,i)$ with $O_{ui}=0$, the
uncorrected DR propagation $\tW\tW^\top\tW$ has conditional bias
\[
C_{ui}\hat Y_{ui}\,(R_u+K_i)-C_{ui}^3\hat Y_{ui}(1-\hat Y_{ui}^2),\qquad
K_i=\sum_{v\neq u}C_{vi}^2\,\tfrac{1-p_{vi}}{p_{vi}}\,\varepsilon_{vi}^2,
\]
and $R_u$ is the analogous sum over the user's own edges. The within-user
ranking bias $C_{ui}\hat Y_{ui}K_i$ satisfies
$K_i\ge(1/\max_vp_{vi}-1)\sum_{v\neq u}C_{vi}^2\varepsilon_{vi}^2$. It therefore
diverges as the item's propensities vanish, and it is $\Omega(1/\tau)$ under
clipping. DRUP's bias is zero.
\end{theorem}

\begin{theorem}[Variance, clipping and concentration]\label{thm:var}
Let $\Lambda_e(u,i)$ be the $C$-weighted mass of the walks from $u$ to $i$
through edge $e$, and let $B=\max(1,1/\tau)$. Then
(a) $\Var(W_e)\le\varepsilon_e^2/\tau$;
(b) $\Var(\hat T_{ui})\le\frac{(1+1/\tau)^2}{\tau}\sum_e\varepsilon_e^2\Lambda_e(u,i)^2$;
(c) $|\E\hat T_{ui}-T^*_{ui}|\le3\bar\delta(1+\bar\delta)^2\sum_{j,v}C_{uj}C_{vj}C_{vi}$ with
$\bar\delta=\max_e(1-p_e/\tau)_+\varepsilon_e$;
(d) $\Pr(|\hat T_{ui}-\E\hat T_{ui}|\ge t)\le2\exp(-2t^2/\sum_ec_e^2)$ with
$c_e=\varepsilon_eB^2\Lambda_e(u,i)/\tau$.
\end{theorem}

The variance scales with the imputation error $\varepsilon^2$, not with $Y^2$
as for IPS, and the clip $\tau$ trades variance against a clipping bias that
vanishes when $\hat Y$ is accurate.

\begin{definition}[Exposure elasticity]
For an intervention $do(p_{\cdot i}\leftarrow\pi p_{\cdot i})$ on item $i$'s
exposure mechanism, $\eta_i=\partial\log\E[s_{ui}]/\partial\log\pi$.
\end{definition}

\begin{theorem}[Exposure invariance: causal item fairness]\label{thm:fair}
Under Assumptions~\ref{as:unconf}--\ref{as:fixed} and $\pi p\ge\tau$, DRUP has
$\eta_i=0$, so items with equal potential outcomes receive equal expected
scores whatever their exposure history. Logged-graph propagation has
$\eta_i\ge1$ (popularity amplification). Uncorrected DR propagation has
$\eta_i<0$ whenever $\hat Y_{ui}K_i>0$ (over-correction).
\end{theorem}

\begin{proposition}[Affine in the user's own data]\label{prop:affine}
Fix $u$ and a candidate $i$ with $O_{ui}=0$, and let $G^{(-u)}$ be $G$ without
$u$'s own contribution. Then
$s_{ui}=\mathrm{const}+\sum_{j\neq i}\big(c_{uj}G^{(-u)}_{ji}+c_{ui}w_{ui}c_{uj}^2\big)w_{uj}$.
Hence un-logging interaction $j$ changes $s_{ui}$ by exactly
$\phi_j=a_j(w_{uj}-\hat Y_{uj})$. Removal effects add up over any subset, the
$\phi_j$ are the exact Shapley values, and completeness holds exactly.
\end{proposition}

\begin{theorem}[Optimal minimal counterfactual explanation]\label{thm:cf}
For candidates $i,k$ with $s_{ui}>s_{uk}$, sorting $g_j=\phi^{(i)}_j-\phi^{(k)}_j$ in
decreasing order and removing interactions until the cumulative sum exceeds the
margin yields a minimum-cardinality set whose removal makes $k$ outrank $i$.
If the positive $g_j$ do not sum above the margin, no such set exists. The
cost is $O(n+d_u\log d_u)$.
\end{theorem}

\begin{theorem}[Certified robustness to injected profiles]\label{thm:robust}
With frozen nuisances, one injected user $v$ changes $u$'s 3-hop score of item
$k$ by $\Delta_k=x_k(D+a_kc_{vk})-a_kx_k^2$, where $a=\tilde w_u$, $x=c_v\odot w_v$
and $D=\langle a,x\rangle$. If $v$ logs at most $L$ interactions with arbitrary
items, labels and propensities $\ge\tau$, then $D$ and $x_k$ lie in closed-form
intervals. $\Delta_k$ is linear in $D$ and quadratic in $x_k$, so its extremes
$[\underline\Delta_k,\overline\Delta_k]$ over this box are exact. Item $t$
provably stays out of $u$'s top-$K$ against any $F$ such users if
$s_{ut}+F\overline\Delta_t$ is below the $(K{+}1)$-th largest value of
$s_{uk}+F\underline\Delta_k$. With unclipped IPS, the influence of a single
interaction is unbounded.
\end{theorem}

\begin{theorem}[Joint differential privacy]\label{thm:dp}
Scale each user row so that $\|\tilde w_v\|_2\le R$ and
$\|c_v^{\odot2}\odot w_v\|_2\le R^2$. Releasing $G$ plus symmetric Gaussian noise with
$\sigma=2R^2\sqrt{2\ln(1.25/\delta)}/\epsilon$ is $(\epsilon,\delta)$-DP under adding
or removing one user. Scoring by~\eqref{eq:local} uses only the release and the
user's own row, so the recommender is $(\epsilon,\delta)$-jointly
DP~\cite{hsu2014matchings}. Post-processing such as low-rank denoising is free.
\end{theorem}

%======================================================================
\section{Experiments}
%======================================================================

\subsection{Setup}
\paragraph{Data.}
\emph{Coat}~\cite{schnabel2016} has 290 users and 300 items. Training
consists of 6{,}960 self-selected (MNAR) ratings. Testing uses 16 ratings per
user on random items (MAR), and we rank only items the user never rated. The
dataset ships propensities; we also fit our own exposure model, with the same
conclusions. \emph{KuaiRec}~\cite{gao2022kuairec} provides the big matrix
(7{,}176 users, 10{,}728 items, 10.3M logged views) as the MNAR log and a fully
observed small matrix (1{,}411 users, 3{,}327 items, 99.6\% dense) as the test
set. We rank all never-logged small-matrix items. Positives are ratings $\ge4$
(Coat) and watch ratio $\ge2$ (KuaiRec).

\paragraph{Protocol.}
Hyper-parameters are chosen on a validation part of the unbiased data: 30\% of
each user's test items for Coat, 30\% of the test users for KuaiRec. We report
the held-out remainder, averaged over 10 (Coat) or 5 (KuaiRec) random splits.
Trained baselines are popularity, imputation only, and MF / IPS-MF / DR-MF /
LightGCN / NAVIP-LightGCN / DR-LightGCN with 64-dimensional embeddings,
pointwise losses and early stopping on the validation part. On Coat they use
10 splits and a grid over learning rate and weight decay. On KuaiRec, because
of our CPU-only budget, they use 3 splits, one learning rate
($10^{-2}$), one weight decay ($10^{-4}$) and at most 15 epochs, so their
KuaiRec numbers should be read as lower bounds. The training-free propagation
family shares one scoring form and differs only in its operator: the logged
graph (linear LightGCN), IPS adjacency (NAVIP-style), IPS with the walk
correction (WC), DR adjacency, and DRUP.

\subsection{Monte-Carlo verification}
Table~\ref{tab:mc} uses a synthetic MNAR model with known propensities,
20{,}000 replications for bias and 3{,}000 for elasticity. The uncorrected IPS
and DR operators have relative biases of 2.4 and 1.9. With the walk correction,
bias disappears: the fraction of entries with $|z|>3$ matches the 0.27\% expected
by chance, and RMSE drops about fivefold. Elasticity follows
Theorem~\ref{thm:fair}: $+1.78$ for the logged graph, $-0.19$ for DR adjacency,
and $\approx0$ for DRUP.

\begin{table}[t]\centering\small
\caption{Monte-Carlo check of Theorems~\ref{thm:unbiased}, \ref{thm:lower} and \ref{thm:fair}.}\label{tab:mc}
\begin{tabular}{lcccc}
\toprule
Estimator & rel.\,|bias| & frac.\ $|z|>3$ & rel.\,RMSE & exposure elasticity $\eta$\\
\midrule
Obs & 0.992 & 1.0000 & 1.02 & +1.777\\
IPS & 2.403 & 0.4042 & 21.01 & +0.007\\
IPS+WC & 0.028 & 0.0025 & 4.93 & +0.007\\
DR & 1.903 & 0.6292 & 20.94 & -0.187\\
DRUP & 0.020 & 0.0008 & 3.97 & +0.005\\
\bottomrule
\end{tabular}

\end{table}

\subsection{Accuracy}
\begin{table}[t]\centering\small
\caption{Accuracy on unbiased test data (mean$\pm$std over splits). Best in bold, second underlined.}\label{tab:acc}
\begin{adjustbox}{max width=\textwidth}
\begin{tabular}{llcccc}
\toprule
 & Method & Coat nDCG@5 & Yahoo!\,R3 nDCG@5 & KuaiRand nDCG@10 & KuaiRec nDCG@20\\
\midrule
\multicolumn{6}{l}{\emph{Non-personalised}}\\
 & Popularity & 0.5367$\pm$0.0095 & 0.5331$\pm$0.0059 & -- & 0.5766$\pm$0.0023\\
 & Imputation only & 0.5532$\pm$0.0097 & 0.5162$\pm$0.0044 & -- & 0.6194$\pm$0.0018\\
\multicolumn{6}{l}{\emph{Trained, pointwise loss}}\\
 & MF & 0.4780$\pm$0.0168 & 0.5076$\pm$0.0060 & -- & 0.6227$\pm$0.0029\\
 & IPS-MF & 0.4713$\pm$0.0160 & 0.5098$\pm$0.0048 & -- & 0.6258$\pm$0.0025\\
 & DR-MF & 0.5237$\pm$0.0111 & 0.5116$\pm$0.0020 & -- & 0.6191$\pm$0.0020\\
 & DR-JL & 0.4811$\pm$0.0164 & -- & -- & --\\
 & MRDR & 0.4869$\pm$0.0187 & -- & -- & --\\
 & iALS & \textbf{0.5812$\pm$0.0101} & -- & -- & --\\
 & LightGCN & 0.5389$\pm$0.0107 & 0.5935$\pm$0.0059 & -- & 0.6277$\pm$0.0027\\
 & DR-LightGCN & 0.5724$\pm$0.0205 & 0.5440$\pm$0.0029 & -- & 0.6246$\pm$0.0016\\
\multicolumn{6}{l}{\emph{Trained, pairwise or sampled loss}}\\
 & MF (BPR) & 0.5174$\pm$0.0206 & 0.6460$\pm$0.0041 & -- & 0.5913$\pm$0.0028\\
 & PDA & 0.5181$\pm$0.0138 & 0.6516$\pm$0.0052 & -- & 0.5822$\pm$0.0020\\
 & MACR & 0.5418$\pm$0.0134 & -- & -- & --\\
 & LightGCN (BPR) & 0.5683$\pm$0.0132 & 0.6693$\pm$0.0045 & -- & 0.6086$\pm$0.0028\\
 & r-AdjNorm & 0.5722$\pm$0.0134 & \textbf{0.6713$\pm$0.0050} & -- & 0.6086$\pm$0.0031\\
 & NAVIP & 0.5664$\pm$0.0163 & \underline{0.6694$\pm$0.0059} & -- & 0.6183$\pm$0.0024\\
 & SimGCL & 0.5528$\pm$0.0141 & -- & -- & --\\
\multicolumn{6}{l}{\emph{Training-free, logged graph}}\\
 & linear LightGCN / r-AdjNorm & 0.5571$\pm$0.0100 & 0.6585$\pm$0.0058 & -- & 0.6036$\pm$0.0028\\
 & EASE & 0.5401$\pm$0.0138 & 0.6478$\pm$0.0061 & -- & 0.5213$\pm$0.0032\\
 & GF-CF & 0.5513$\pm$0.0123 & 0.6552$\pm$0.0039 & -- & 0.6141$\pm$0.0025\\
 & BSPM & \underline{0.5727$\pm$0.0090} & 0.6507$\pm$0.0047 & -- & --\\
\multicolumn{6}{l}{\emph{Training-free, debiased graph}}\\
 & IPS adjacency (NAVIP-style) & 0.5510$\pm$0.0224 & 0.6567$\pm$0.0063 & -- & 0.6132$\pm$0.0030\\
 & IPS + walk correction & 0.5510$\pm$0.0224 & 0.6567$\pm$0.0063 & -- & 0.6132$\pm$0.0030\\
 & EASE on DR graph & 0.5564$\pm$0.0141 & 0.6072$\pm$0.0033 & -- & 0.6026$\pm$0.0022\\
 & GF-CF on DR graph & 0.5608$\pm$0.0083 & 0.6312$\pm$0.0051 & -- & 0.6371$\pm$0.0021\\
 & BSPM on DR graph & 0.5624$\pm$0.0112 & 0.6262$\pm$0.0071 & -- & --\\
 & DR adjacency & 0.5459$\pm$0.0244 & 0.6567$\pm$0.0063 & -- & \underline{0.6378$\pm$0.0020}\\
 & DRUP & 0.5467$\pm$0.0235 & 0.6567$\pm$0.0063 & -- & 0.6378$\pm$0.0019\\
 & DR adjacency, 5 hops & 0.5524$\pm$0.0245 & -- & -- & \textbf{0.6379$\pm$0.0011}\\
 & DRUP, 5 hops & 0.5451$\pm$0.0204 & -- & -- & 0.6376$\pm$0.0011\\
\multicolumn{6}{l}{\emph{Training-free, sample split (Assumption 2 holds)}}\\
 & DR adjacency, sample split & 0.5604$\pm$0.0132 & 0.5934$\pm$0.0063 & -- & 0.6338$\pm$0.0013\\
 & DRUP-split & 0.5578$\pm$0.0144 & 0.5925$\pm$0.0051 & -- & 0.6337$\pm$0.0013\\
\bottomrule
\end{tabular}

\end{adjustbox}
\end{table}

Table~\ref{tab:acc} reports accuracy and Table~\ref{tab:sig} the paired
$t$-tests over shared splits. DRUP needs no training. On KuaiRec it has the
highest nDCG@20 (0.6326). It is significantly better than the imputation
model, logged-graph and IPS propagation, DR-MF and DR-LightGCN, and it is
ahead of LightGCN, MF and NAVIP by margins that are not significant with three
splits. On Coat it significantly outperforms MF, IPS-MF, DR-MF, LightGCN and
NAVIP and ties with DR-LightGCN. It is, however, significantly \emph{below}
the imputation-only model ($-0.014$) and the uncorrected DR adjacency
($-0.006$). With 290 users and 16 random ratings each, the additive imputation
already captures most of the recoverable signal. The walk correction therefore
does not buy top-$K$ accuracy. It buys unbiasedness, exposure invariance and
certifiability. As predicted by Theorem~\ref{thm:lower}, IPS+WC coincides with
IPS on unexposed candidates at three hops, because the factor $W_{ui}=0$ hides
the IPS bias. We do not claim a large accuracy gain. The claim is that DRUP is
accuracy-competitive with trained models while adding the guarantees of
Section~5.

\begin{table}[t]\centering\small
\caption{Paired $t$-tests of DRUP against each baseline (nDCG@5 on Coat, nDCG@20 on KuaiRec).}\label{tab:sig}
\begin{tabular}{lcccccc}
\toprule
DRUP against & \multicolumn{2}{c}{Coat} & \multicolumn{2}{c}{Yahoo!\,R3} & \multicolumn{2}{c}{KuaiRec}\\
 & $\Delta$ [95\% CI] & $p_{\mathrm{Holm}}$ & $\Delta$ [95\% CI] & $p_{\mathrm{Holm}}$ & $\Delta$ [95\% CI] & $p_{\mathrm{Holm}}$\\
\midrule
DR adjacency (effect of the correction) & +0.0020 [-0.003, +0.007] & 1.000 & -- & -- & -- & --\\
DR adjacency, 5 hops (effect of the correction) (5 hops) & -- & -- & -- & -- & -- & --\\
linear LightGCN on the logged graph & -0.0097 [-0.049, +0.029] & 1.000 & -- & -- & -- & --\\
imputation only & -0.0057 [-0.016, +0.004] & 1.000 & -- & -- & -- & --\\
GF-CF & -0.0043 [-0.043, +0.034] & 1.000 & -- & -- & -- & --\\
EASE & +0.0061 [-0.033, +0.045] & 1.000 & -- & -- & -- & --\\
GF-CF on DR graph & -0.0102 [-0.017, -0.004] & 0.058 & -- & -- & -- & --\\
EASE on DR graph & -0.0074 [-0.021, +0.006] & 1.000 & -- & -- & -- & --\\
LightGCN (BPR) & -0.0248 [-0.058, +0.008] & 1.000 & -- & -- & -- & --\\
r-AdjNorm & -0.0294 [-0.063, +0.004] & 1.000 & -- & -- & -- & --\\
NAVIP & -0.0212 [-0.054, +0.012] & 1.000 & -- & -- & -- & --\\
MF (BPR) & +0.0372 [+0.007, +0.067] & 0.287 & -- & -- & -- & --\\
DR-MF & +0.0235 [+0.002, +0.045] & 0.627 & -- & -- & -- & --\\
PDA & +0.0277 [-0.004, +0.060] & 1.000 & -- & -- & -- & --\\
DR-LightGCN & +0.0004 [-0.021, +0.022] & 1.000 & -- & -- & -- & --\\
\bottomrule
\end{tabular}

\end{table}

\subsection{Fairness}
\begin{table}[t]\centering\small
\caption{Fairness. ECB is the partial correlation between an item's mean rank and its log exposure, controlling for its true like rate. The user gap is the nDCG gap between user groups (gender on Coat, activity on KuaiRec). The last column is the rank shift of items whose exposures were thinned by a real $do(p\leftarrow p/2)$; 0 means exposure invariance.}\label{tab:fair}
\begin{adjustbox}{max width=\textwidth}
\begin{tabular}{llcccccc}
\toprule
Data & Method & nDCG & ECB$\downarrow$ & PRU$\downarrow$ & Gini$\downarrow$ & user gap$\downarrow$ & rank shift under $do(p/2)$\\
\midrule
Coat & Obs & 0.4917 & 0.769 & 0.394 & 0.340 & 0.1123 & -0.0585$\pm$0.0073\\
Coat & IPS & 0.4864 & 0.765 & 0.392 & 0.337 & 0.1189 & -0.0594$\pm$0.0066\\
Coat & DR & 0.4770 & 0.744 & 0.763 & 0.655 & 0.0795 & -0.0027$\pm$0.0122\\
Coat & DRUP & 0.4803 & 0.739 & 0.758 & 0.655 & 0.0793 & +0.0006$\pm$0.0110\\
\bottomrule
\end{tabular}

\end{adjustbox}
\end{table}

Table~\ref{tab:fair} tests Theorem~\ref{thm:fair} on real data. We thin the
logged exposures of a random half of the items with probability $1/2$, which
is an intervention with a known propensity change. Logged-graph and
(clipped) IPS propagation demote the thinned items, i.e., they punish
under-exposure. DRUP's rank shift is statistically indistinguishable from zero.
The NDCG gap between user groups is also smaller for the DR-based operators.
\emph{Limitation:} exposure invariance is a causal guarantee, not a
distributional one. When true quality is concentrated, DRUP's top-$K$ lists
are also concentrated, and its Gini and PRU are higher than those of the
logged graph on Coat. The accuracy--concentration frontier in
Table~\ref{tab:frontier} makes this trade-off explicit. Combining DRUP with an
exposure-constrained re-ranker is left for future work.

\begin{table}[t]\centering\small
\caption{Best nDCG reachable by each operator under a cap on the Gini coefficient of top-$K$ exposure (over the whole hyper-parameter grid).}\label{tab:frontier}
\begin{tabular}{llccccc}
\toprule
Data & Method & Gini$\le$0.3 & $\le$0.4 & $\le$0.5 & $\le$0.6 & any\\
\midrule
Coat & Obs & -- & 0.4917 & 0.4917 & 0.4917 & 0.4917\\
Coat & IPS & -- & 0.4864 & 0.4864 & 0.4864 & 0.4864\\
Coat & DR & -- & 0.4090 & 0.4487 & 0.4733 & 0.4853\\
Coat & DRUP & -- & 0.3786 & 0.4416 & 0.4848 & 0.4902\\
\bottomrule
\end{tabular}

\end{table}

\subsection{Transparency}
\begin{table}[t]\centering\small
\caption{Exact explanations (Prop.~\ref{prop:affine}, Thm.~\ref{thm:cf}) on 200 random test users.}\label{tab:explain}
\begin{tabular}{lcccc}
\toprule
Data & max.\ rel.\ error & top-1 rank after 5 removals (attr.\ / random / similarity) & no CF explanation & median CF size\\
\midrule
Coat & 6.1e-16 & 0.04 / 0.02 / 0.01 & 67\% & 0.0\\
Yahoo!\,R3 & 8.7e-16 & 1.65 / 0.27 / 1.40 & 16\% & 1.0\\
\bottomrule
\end{tabular}

\end{table}

Attributions recomputed against full re-propagation are exact to machine
precision, both for completeness and for additivity over random subsets
(Table~\ref{tab:explain}). Minimal counterfactual explanations are small, and
where none exists the explanation is itself informative. On Coat the top
recommendation of most users cannot be overturned by removing any subset of
their own interactions: it is driven by population-level evidence, which DRUP
can certify.

\subsection{Accountability: certified robustness}
\begin{table}[t]\centering\small
\caption{Injection attack: $F$ fake users log a low-popularity target item and the most popular items. Left: realised hit rate of the target in real users' top-$K$. Right: fraction of real users for whom Theorem~\ref{thm:robust} certifies that the target cannot enter the top-$K$ against \emph{any} $F$ profiles of that size. No certificate was violated.}\label{tab:attack}
\begin{adjustbox}{max width=\textwidth}
\begin{tabular}{lcccccccccccc}
\toprule
 & \multicolumn{6}{c}{hit rate of target in top-$K$ $\downarrow$} & \multicolumn{6}{c}{certified fraction $\uparrow$}\\
\midrule
\multicolumn{13}{l}{\textit{Coat} (K=5, profile size 21)}\\
Method & $F$=0 & $F$=1 & $F$=5 & $F$=10 & $F$=20 & $F$=50 & $F$=0 & $F$=1 & $F$=5 & $F$=10 & $F$=20 & $F$=50\\
Obs & 0.004 & 0.004 & 0.009 & 0.108 & 0.444 & 0.661 & 0.996 & 0.000 & 0.000 & 0.000 & 0.000 & 0.000\\
IPS & 0.004 & 0.004 & 0.020 & 0.201 & 0.502 & 0.651 & 0.995 & 0.000 & 0.000 & 0.000 & 0.000 & 0.000\\
DR & 0.000 & 0.000 & 0.000 & 0.000 & 0.000 & 0.655 & 1.000 & 0.000 & 0.000 & 0.000 & 0.000 & 0.000\\
DRUP & 0.000 & 0.000 & 0.000 & 0.000 & 0.382 & 0.973 & 1.000 & 1.000 & 0.211 & 0.000 & 0.000 & 0.000\\
\bottomrule
\end{tabular}

\end{adjustbox}
\end{table}

Table~\ref{tab:attack} shows that clipping plus imputation-based degrees
make DRUP the only operator with non-trivial certificates. The certificate
was never violated by a realised attack, and brute force over all small
profiles confirms that the bound is valid and reaches 62--67\% of the true
worst case. Empirically, DR adjacency with attacker-controlled degrees can be
harder to push at intermediate budgets, because the fake user's large DR
weights inflate its own degree. That protection is uncertifiable, since the
attacker can choose a profile that avoids it.

\subsection{Privacy}
\begin{table}[t]\centering\small
\caption{nDCG of jointly DP DRUP ($\delta=10^{-5}$). The denoising rank is chosen on validation data.}\label{tab:privacy}
\begin{adjustbox}{max width=\textwidth}
\begin{tabular}{llccccccc}
\toprule
Data & rank & $\epsilon=0.5$ & 1 & 2 & 4 & 8 & 16 & $\infty$\\
\midrule
Coat & 8 & 0.381 & 0.385 & 0.400 & 0.399 & 0.398 & 0.449 & 0.533\\
 & 32 & 0.382 & 0.385 & 0.403 & 0.385 & 0.401 & 0.438 & 0.548\\
 & 128 & 0.372 & 0.370 & 0.387 & 0.401 & 0.389 & 0.422 & 0.536\\
 & full & 0.376 & 0.389 & 0.378 & 0.393 & 0.394 & 0.417 & 0.532\\
 & \multicolumn{8}{l}{\footnotesize non-private DRUP (all nuisances) 0.555; popularity 0.527; 29 public users}\\
\midrule
Yahoo!\,R3 & 8 & 0.483 & 0.492 & 0.491 & 0.493 & 0.492 & 0.493 & 0.492\\
 & 32 & 0.481 & 0.490 & 0.491 & 0.493 & 0.492 & 0.493 & 0.492\\
 & 128 & 0.479 & 0.490 & 0.490 & 0.494 & 0.493 & 0.492 & 0.493\\
 & full & 0.473 & 0.488 & 0.491 & 0.494 & 0.494 & 0.493 & 0.492\\
 & \multicolumn{8}{l}{\footnotesize non-private DRUP (all nuisances) 0.652; popularity 0.525; 1540 public users}\\
\midrule
KuaiRec & 8 & 0.175 & 0.179 & 0.173 & 0.427 & 0.603 & 0.618 & 0.634\\
 & 32 & 0.174 & 0.175 & 0.154 & 0.412 & 0.598 & 0.618 & 0.634\\
 & 128 & 0.160 & 0.173 & 0.172 & 0.365 & 0.581 & 0.614 & 0.634\\
 & full & 0.182 & 0.202 & 0.223 & 0.280 & 0.377 & 0.492 & 0.634\\
 & \multicolumn{8}{l}{\footnotesize non-private DRUP (all nuisances) 0.637; popularity 0.577; 717 public users}\\
\bottomrule
\end{tabular}

\end{adjustbox}
\end{table}

Table~\ref{tab:privacy} shows the privacy--utility curve. Coat has only 290
users, so meaningful privacy costs accuracy there. KuaiRec, with 25 times more
users, supports much smaller $\epsilon$.

%======================================================================
\section{Discussion and limitations}
%======================================================================
DRUP relies on unconfoundedness and on nuisances that are fixed or cross-fitted,
the standard assumptions of IPS/DR estimation. The exact correction is derived
for up to three hops (and for the item Gram operator). Higher orders need the
same idempotent rule, but the number of repetition patterns grows. The fairness
guarantee concerns causal exposure invariance, not parity of outcomes. The
privacy analysis treats the low-dimensional nuisances as public. Beyond
accuracy, DRUP turns a graph recommender into an object whose bias, fairness,
explanations, robustness and privacy can be stated and checked.

\bibliographystyle{abbrvnat}
\bibliography{refs}

\appendix
\section{Proofs}\label{app:proofs}

\paragraph{Theorem~\ref{thm:unbiased}.}
After the correction every walk contributes $\prod_{e\in\mathrm{distinct}}C_e^{k_e}W_e$.
The $W_e$ are independent (Assumption~\ref{as:unconf}, fixed nuisances), so
$\E\prod W_e=\prod\E W_e$ with $\E W_e=\hat Y_e+p_e(Y_e-\hat Y_e)/\bar p_e=Y_e+\delta_e$.
The target monomial is $\prod C_e^{k_e}Y_e^{k_e}=\prod C_e^{k_e}Y_e$ because
$Y_e\in\{0,1\}$. \qed

\paragraph{Theorem~\ref{thm:lower}.}
The naive operator contains, for $v=u$, $j\ne i$, the terms
$\tW_{uj}^2\tW_{ui}$; for $j=i$, $v\ne u$, the terms $\tW_{ui}\tW_{vi}^2$; and
$\tW_{ui}^3$. Conditionally on $O_{ui}=0$ we have $W_{ui}=\hat Y_{ui}$. By
independence, $\E[W_e^2]=Y_e+\Var(W_e)$ with
$\Var(W_e)=\frac{1-p_e}{p_e}(Y_e-\hat Y_e)^2$, while the corrected replacement has
mean $Y_e$. Summing the excesses gives the formula. The lower bound uses
$(1-p)/p\ge1/\max p-1$. \qed

\paragraph{Theorem~\ref{thm:var}.}
(a) $\Var(W_e)=p_e(1-p_e)\varepsilon_e^2/\bar p_e^2$. If $p_e\ge\tau$ this is at
most $(1/\tau-1)\varepsilon_e^2$; otherwise it is at most $p_e\varepsilon_e^2/\tau^2<\varepsilon_e^2/\tau$.
(b) $\hat T_{ui}$ is multilinear in independent variables and therefore affine
in each $W_e$. The Efron--Stein inequality~\cite{efron1981jackknife} gives
$\Var(\hat T)\le\sum_e\E[(\partial_e\hat T)^2]\Var(W_e)$. The derivative
$\partial_e\hat T$ is a sum, over the walks through $e$, of $C$-weights times at
most two other independent $W_f$ with $\|W_f\|_2^2\le1+1/\tau$. Minkowski's
inequality then bounds $\|\partial_e\hat T\|_2\le(1+1/\tau)\Lambda_e$.
(c) By Theorem~\ref{thm:unbiased}, the bias is at most
$\sum|\kappa_{\mathrm{walk}}|[\prod(Y_f+|\delta_f|)-\prod Y_f]\le((1+\bar\delta)^3-1)M_{ui}$,
and $|\delta_e|=(1-p_e/\tau)_+\varepsilon_e$.
(d) Changing $O_e$ moves $W_e$ by at most $\varepsilon_e/\tau$, and
$|\partial_e\hat T|\le B^2\Lambda_e$. Apply McDiarmid's inequality~\cite{mcdiarmid1989}. \qed

\paragraph{Theorem~\ref{thm:fair}.}
DRUP: by Theorem~\ref{thm:unbiased}, $\E[s]=F^*$, which does not depend on $\pi$.
Logged graph: $\E[O_eY_e]=p_eY_e$, so each 3-hop monomial containing
$k\ge1$ edges incident to $i$ scales as $\pi^k$. Uncorrected DR: by
Theorem~\ref{thm:lower}, $\E[s]=F^*+C_{ui}\hat Y_{ui}K_i(\pi)+\dots$, and
$K_i(\pi)$ is strictly decreasing in $\pi$. \qed

\paragraph{Proposition~\ref{prop:affine} and Theorem~\ref{thm:cf}.}
Substitute $G=G^{(-u)}+\tilde w_u^\top\tilde w_u-\diag(\tilde w_u^2)+\diag(c_u^2w_u)$
into~\eqref{eq:local}. The terms $\tilde w_{ui}\sum_j\tilde w_{uj}^2$ cancel
against the local correction, leaving $\tilde w_{ui}\sum_jc_{uj}^2w_{uj}$, and
every remaining term involving $w_{uj}$ ($j\neq i$) is linear. Additive games
have Shapley values equal to their individual effects. For
Theorem~\ref{thm:cf}, the margin after removing $S$ is
$s_{ui}-s_{uk}-\sum_{j\in S}g_j$, so minimising $|S|$ under
$\sum_Sg_j>$ margin selects the largest $g_j$. \qed

\paragraph{Theorem~\ref{thm:robust}.}
Adding row $v$ adds $\mathrm{contrib}_v$ to $G$ and leaves $u$'s local terms
unchanged. Then $\Delta=\tilde w_u\,\mathrm{contrib}_v=(a\cdot x)x-a\odot x^2+a\odot c_v\odot x$.
Logged DR entries satisfy
$w\in[\hat Y-\hat Y/\tau,\hat Y+(1-\hat Y)/\tau]$ for any label and any $p\ge\tau$,
and unlogged ones equal $\hat Y$. $D$ is maximised (minimised) by assigning the
$L$ slots to the largest positive (negative) deviations. For fixed $x_k$ the
function is linear in $D$, so it attains its extremes at the ends of the $D$
interval. For fixed $D$ it is a quadratic in $x_k$, whose extremes over an interval
lie at the ends or at the clamped vertex. Effects add over fake users. For
unclipped IPS, $|W_e|=1/p_e$ is unbounded. \qed

\paragraph{Theorem~\ref{thm:dp}.}
$\|\mathrm{contrib}_v\|_F\le\|\tilde w_v\tilde w_v^\top-\diag(\tilde w_v^2)\|_F+\|c_v^2w_v\|_2\le R^2+R^2$,
so the $\ell_2$-sensitivity of the upper triangle is at most $2R^2$, and the
Gaussian mechanism~\cite{dwork2014algorithmic} applies. Each user's output is a
function of the DP release and their own data, which gives joint DP by the
billboard lemma~\cite{hsu2014matchings}. \qed

\end{document}
